\documentclass[11pt]{article}
\usepackage[T1]{fontenc}
\usepackage{amsmath,amssymb,amsthm}
\usepackage[margin=1.1in]{geometry}
\usepackage{microtype}
\usepackage[hidelinks]{hyperref}
\setlength{\emergencystretch}{3em}

\newtheoremstyle{erratumstyle}{}{}{}{}{\bfseries}{.}{ }{#1 #2 (#3)}
\theoremstyle{erratumstyle}
\newtheorem{erratum}{Erratum}

\title{Errata for\\
  \emph{A Time and Space Optimal Stable Population Protocol\\
  Solving Exact Majority}\\[2mm]
  {\large Doty, Eftekhari, G\k{a}sieniec, Severson, Stachowiak, Uzna\'nski}\\[4mm]
  {\large found while formalizing the paper in Lean 4 / Mathlib}}
\author{}
\date{}

\begin{document}
\maketitle

\begin{abstract}
We formalized Theorem~3.1 of the paper in Lean~4 with Mathlib: the protocol
\textsc{Nonuniform Majority} uses $O(\log n)$ states, stably computes
majority, and stabilizes in $O(\log n)$ parallel time with probability
$1-O(1/n^2)$ and in expectation. Theorem~3.1 holds. Along the way we found
the items below. Two printed statements are false as stated: the first case
of Lemma~5.3 for large gaps, and the lower bound of Theorem~6.9 for small
clock fractions, which comes from a wrong conversion factor. Neither affects
Theorem~3.1. The first needs one extra step in the expected-time argument;
the second leaves the choice $k=45$ in Theorem~6.12 without a proof, and a
larger constant $k$ restores it. The other items are proof gaps with true
statements, and typographical or pseudocode slips whose intended meaning is
clear from the text.

All locations refer to arXiv:2106.10201v2; page numbers are those of its PDF.
Classification: (E) a printed statement is false; (P) the statement is true
but the printed proof does not establish it; (T) a typographical or pseudocode
slip with unambiguous intended meaning.
\end{abstract}

\section*{Notation}

As in the paper, $n$ is the population size, $g$ the initial gap, $M$ the set
of Main agents, $c=|C|/n$ the fraction of Clock agents, $m=|M|/n$, and $k$ the
number of minutes per hour. Parallel time is the number of interactions
divided by $n$.

%--------------------------------------------------------------------
\begin{erratum}[Lemma 5.3, large gaps (E)]
\textbf{Location.} Lemma~5.3 and its proof (p.~26); the description of the
Phase~2 Init (p.~18); the expected-time paragraph of the proof of
Theorem~3.1 (p.~46).

\medskip
\textbf{As printed.} ``If the initial gap $|g|\ge 0.025|M|$, then we stabilize
to the correct output in Phase~2'' with probability $1-O(1/n^2)$. The proof
treats $|g|>0.5|M|$ by noting $|\mu|\ge1$, ``so all remaining biased agents
have the majority opinion, and we will stabilize in Phase~2 to the correct
majority output.'' On p.~18 the check $|\mathrm{bias}|>1$ in the Phase~2 Init
``can only happen with low probability.''

\medskip
\textbf{Issue.} Phases~0 and~1 conserve the sum of the biases of Main and
RoleMCR agents, which equals $g$; a RoleMCR agent that reaches Phase~1 sends
the population to Phase~10. Main agents leave Phase~1 only through lines
1--4 of the main loop, which execute the Phase~2 Init, and no Phase~2 rule
changes bias. If $|g|>|M|$, consider the last Main agent to leave Phase~1.
Every other Main agent left with $|\mathrm{bias}|\le1$ or has already
triggered the error, so this agent holds bias of absolute value at least
$|g|-(|M|-1)\ge2$. Its Phase~2 Init sends it to Phase~10 (reading the main
loop as in Erratum~9), and the epidemic brings the whole population to the
backup. Thus for every input with $|g|>|M|$ the protocol reaches Phase~10
with probability~$1$ once Phase~1 has started. Since $|M|\le 2n/3$ at all
times (from the invariant $s_f+2s_t=m_f+2m_t$ in the proof of Lemma~5.2; the
lemma states the bound under the condition that Phase~1 initializes without
error), this includes every input with $|g|>2n/3$, in particular the
unanimous input $g=n$. The range $0.5|M|<|g|\le|M|$ has a related step to
check: rounding gives $\mu=1$, and the cited convergence to
$\{\mu-1,\mu,\mu+1\}=\{0,1,2\}$ does not by itself exclude an agent with
bias~$2$.

\medskip
\textbf{Correction.} The output remains correct, since Phase~10 is a stable
backup. For the time bound, the expected-time paragraph on p.~46 charges
Phase~10 only on an event of probability $O(1/n^2)$ and uses the
$O(n\log n)$ bound of Lemma~7.7 there. For large gaps the backup is fast.
Suppose some agent is in Phase~10 and $|g|\ge n/120$ (which holds whenever
$|g|\ge0.025|M|$, since $|M|\ge n/3$). The phase epidemic to Phase~10
completes in $O(\log n)$ parallel time. In an all-Phase-10 configuration the
number of active majority agents exceeds the number of active minority
agents by $|g|$, and the backup stabilizes after three monotone stages
(eliminate the active minority, then the active ties, then correct the
passive outputs), each with useful-partner probability at least
$2|g|/(n(n-1))$ per target agent. This gives expected stabilization within
$180(1+\ln n)$ parallel time from an all-Phase-10 configuration, and
$O(\log n)$ parallel time with probability $1-O(n^{-3})$. With this step,
Theorem~3.1 holds for all inputs.
\end{erratum}

%--------------------------------------------------------------------
\begin{erratum}[Theorem 6.9, conversion to the whole population (E)]
\textbf{Location.} Theorem~6.9 and its proof (p.~33); its use with $k=45$ in
the proof of Theorem~6.12 (p.~35) and in the proof of Lemma~6.14 (p.~36).

\medskip
\textbf{As printed.} The lemmas of \S6.1 assume $|C|=n$. The proof of
Theorem~6.9 then multiplies their parallel-time bounds by $1/c^2$, since two
Clock agents meet with probability about $c^2$, and obtains
$\mathrm{end}_h-\mathrm{start}_h\ge\frac{1}{c^2}[0.45k-3.1]$ and
$\mathrm{start}_{h+1}-\mathrm{start}_h\le\frac{1}{c^2}[2.11k+2.2]$.

\medskip
\textbf{Issue.} The factor $1/c^2$ converts a number of Clock--Clock
interactions into a number of global interactions. The parallel time of the
all-clock lemmas, however, is measured in units of $|C|$ interactions, while
global parallel time is measured in units of $n$. The correct factor is
\[
  \frac{|C|}{n}\cdot\frac{n(n-1)}{|C|(|C|-1)}=\frac{n-1}{|C|-1}\sim\frac1c .
\]
The printed upper bound remains valid (it is weaker than the corrected one,
$\frac1c[2.11k+2.2]$). The lower bound on the length of a synchronous hour
becomes about $\frac1c[0.45k-3.1]$, smaller by a factor of~$c$. Since
$\mathrm{end}_h\le\mathrm{start}_{h+1}$, the printed lower bound exceeds the
corrected upper bound when $c<(0.45k-3.1)/(2.11k+2.2)$ (about $0.18$ for
$k=45$), so as a statement about an arbitrary clock fraction~$c$ it does not
hold.

\medskip
\textbf{Effect.} Theorem~6.12 needs each synchronous hour to last at least
$\frac2c+\frac{47}{m}\le\frac{49}{c}$ (using $c<m$). It uses
$\frac{0.45\cdot45-3.1}{c^2}\ge\frac{17}{c^2}\ge\frac{49}{c}$. With the
corrected factor, $k=45$ gives $\frac{17.15}{c}<\frac{49}{c}$, so the
argument no longer applies; we do not claim that the conclusion fails. The
proof of Lemma~6.14 uses $[1.48(45-2)-2.2]/c^2\ge 17/c^2\ge 51/m$ in the same
way. (The factor $1.48$ there should be $0.45$, the per-minute lower bound
for $p=1$; the printed chain is unaffected, since $0.45\cdot43-2.2=17.15$.)

\medskip
\textbf{Correction.} Take a larger constant number of minutes per hour.
Theorem~6.12 needs $0.45k-3.1\ge49$, that is $k\ge116$, and Lemma~6.14 needs
$0.45k-3.1\ge51$, that is $k\ge121$. With a $5\%$ allowance for the
concentration factors, $k=128$ satisfies both:
$0.95(0.45\cdot128-3.1)=51.775$. This changes only a constant; the state
count stays $O(\log n)$ and the asymptotic statements are unaffected. Our
formal proof uses $k=256$, which leaves more room in the later phases.
\end{erratum}

%--------------------------------------------------------------------
\begin{erratum}[Lemma 6.7, the displayed integral (T)]
\textbf{Location.} Proof of Lemma~6.7 (p.~31), last display.

\medskip
\textbf{As printed.}
\[
  c_{\ge i+1}(t^{0.1}_i+0.45)\lesssim(1+\epsilon)e^{0.9}
  \left(0.01+\int_0^{0.45}(1+s)^2e^{-2s}\,ds\right)\le0.45 .
\]

\medskip
\textbf{Issue.} The preceding text bounds the drip probability by
$(0.1+s)^2$, so the integrand should be $(0.1+s)^2e^{-2s}$. With $(1+s)^2$
the value is about $1.07$. Also, the lemma needs
$c_{\ge i+1}(t^{0.1}_i+0.45)<0.1$, which the bound $0.45$ does not give.
(The sum before the integral counts agents rather than fractions; a factor
$1/n$ is missing there.)

\medskip
\textbf{Correction.} With the intended integrand,
\[
  e^{0.9}\left(0.01+\int_0^{0.45}(0.1+s)^2e^{-2s}\,ds\right)
  =0.315\,e^{0.9}-0.67625\approx0.0985<0.1,
\]
so the lemma holds for $\epsilon$ small enough.
\end{erratum}

%--------------------------------------------------------------------
\begin{erratum}[Section 6, asynchronous entry into Phase 3 (P)]
\textbf{Location.} \S3.4 (p.~11): ``To simplify the analysis, we formally
start our arguments for each phase assuming each agent is in the current
phase, although technically $\Theta(\log n)$ time will pass between the time
the first agent enters phase~$i$ and the last agent does''; its use in \S6
(the clock lemmas of \S6.1 and the induction on synchronous hours in \S6.2).

\medskip
\textbf{Issue.} The paper states this simplification explicitly. For
Phase~3 it needs an argument. The analysis of \S6 starts from a
configuration in which all Clock agents are at minute~$0$ and the whole
population is in Phase~3 (the base case $c_{\ge0}(0)=1$ of Theorem~6.5,
$\Phi(0)=0$ in Lemma~6.10, the hour-$0$ case of the induction). In the
protocol, agents enter Phase~3 by epidemic over $\Theta(\log n)$ parallel
time, so the first Clock agents drip and drag O agents forward while other
agents are still in Phase~2 (or Phase~1). By the time the last agents enter,
the clock can be $\Theta(\log n)$ minutes ahead. A configuration in which the
whole population has just entered Phase~3 with all minutes at~$0$ occurs only
with probability $e^{-\Omega(n)}$, so the analysis cannot be restarted there.

\medskip
\textbf{Correction.} The clock argument has to start at the first entry into
Phase~3, a stopping time, using the strong Markov property, and account for
the width of the entry front. We did this in our formal proof. The
conclusions used later (synchronous hours of length $\Theta(1/c)$, and the
output of Theorems~6.1 and~6.2 at the end of Phase~3) hold with changed
constants. The same treatment is needed at every phase boundary where the
paper says ``before the first Clock agent advances to the next phase.''
\end{erratum}

%--------------------------------------------------------------------
\begin{erratum}[Lemma 6.10, time horizon (P)]
\textbf{Location.} Proof of Lemma~6.10 (pp.~33--34).

\medskip
\textbf{As printed.} Azuma's inequality is applied to $\Phi$ over
$n\cdot\mathrm{end}_h$ steps, giving $\exp(-O(n)\delta^2)$, ``since by
Theorem~6.9, we have time $\mathrm{end}_h=O(1)$.''

\medskip
\textbf{Issue.} Time is measured from the start of Phase~3, so
$\mathrm{end}_h$ is $\Theta(h)$, which is $\Theta(\log n)$ for the last
hours. The horizon is then $O(n\log n)$ steps and the bound is
$\exp(-\Omega(n/\log n))$.

\medskip
\textbf{Correction.} This is still a very high probability bound, and the
union over hours goes through.
\end{erratum}

%--------------------------------------------------------------------
\begin{erratum}[Hour of a Clock agent, floor or ceiling (T)]
\textbf{Location.} \S3.4, Phase~3 description (p.~18) and pseudocode
(p.~19); the definition of $a.\mathrm{hour}$ (p.~32); proof of Theorem~6.9
(p.~33).

\medskip
\textbf{Issue.} The text on p.~18 defines
$\mathrm{hour}=\lceil\mathrm{minute}/k\rceil$, while the clock update reaction
uses $\lfloor c.\mathrm{minute}/k\rfloor$, p.~32 defines
$a.\mathrm{hour}=\lfloor a.\mathrm{minute}/k\rfloor$, and the proof of
Theorem~6.9 uses $c.\mathrm{hour}=h\iff hk\le c.\mathrm{minute}<(h+1)k$.

\medskip
\textbf{Correction.} Read the floor throughout. With the floor, hour~$L$
corresponds to the single value $\mathrm{minute}=kL$; this is harmless,
since Clock agents at $\mathrm{minute}=kL$ then wait $c_3\ln n$ on their
counter before Phase~3 ends.
\end{erratum}

%--------------------------------------------------------------------
\begin{erratum}[Phase 7 pseudocode, gap-2 cancel (T)]
\textbf{Location.} Phase~7 pseudocode, lines 7--10 (p.~20).

\medskip
\textbf{Issue.} Executed in order, line~9 sets
$i.\mathrm{exponent}\leftarrow i.\mathrm{exponent}-1$ and line~10 then sets
$j.\mathrm{exponent}\leftarrow i.\mathrm{exponent}-2$ with the new value, so
$j$ ends three below the original exponent of~$i$. The example in the
comment, $+\frac14,-\frac1{16}\to+\frac18,+\frac1{16}$, and conservation of
mass ($\frac14-\frac1{16}=\frac18+\frac1{16}$) require two below.

\medskip
\textbf{Correction.} By the guard on line~7, $j$ already has exponent two
below the original exponent of~$i$, so line~10 can be deleted (or executed
before line~9).
\end{erratum}

%--------------------------------------------------------------------
\begin{erratum}[Phase 8 guard (T)]
\textbf{Location.} Phase~8 pseudocode, line~1 (p.~21).

\medskip
\textbf{Issue.} The guard $\{u.\mathrm{bias},v.\mathrm{bias}\}=\{-1,+1\}$
read literally requires both agents at exponent~$0$. The description of
Phase~8 (p.~20) concerns ``opposite-opinion agents of any two exponents,''
and line~2 compares exponents.

\medskip
\textbf{Correction.} The guard should read
$\{u.\mathrm{opinion},v.\mathrm{opinion}\}=\{-1,+1\}$.
\end{erratum}

%--------------------------------------------------------------------
\begin{erratum}[Main loop, phase update after an Init (T)]
\textbf{Location.} \textsc{Nonuniform Majority}, lines 1--4 (p.~12).

\medskip
\textbf{Issue.} If one of the Inits executed in line~3 moves agent~$i$ to
Phase~10 (the error checks in the Inits of Phase~1 and Phase~2), the
assignment $i.\mathrm{phase}\leftarrow j.\mathrm{phase}$ in line~4 moves it
back down, and the loop runs the later Inits. Elsewhere the paper treats the
phase as nondecreasing (p.~11). Erratum~1 uses the intended reading.

\medskip
\textbf{Correction.} Stop the loop once $i$ reaches Phase~10, and set
$i.\mathrm{phase}\leftarrow\max(i.\mathrm{phase},j.\mathrm{phase})$.
\end{erratum}

%--------------------------------------------------------------------
\begin{erratum}[Small slips (T)]
\begin{enumerate}
\item p.~18: the O agents' field is introduced as ``$\mathrm{hour}=L$'',
  while the Init sets $\mathrm{hour}\leftarrow0$; and ``will only participate
  in split reactions with $-\mathrm{exponent}>\mathrm{hour}$'' is reversed
  relative to pseudocode line~13 and Fig.~4c.
\item p.~33: the definition of $\mathrm{end}_h$ uses ``$\le0.001|C|$'' where
  ``reaches $0.001$'' (that is, $\ge$) is meant; in the proof of
  Theorem~6.9, $\mathrm{end}_h=t^{0.001}_{hk}$ should be indexed $(h+1)k$,
  and ``$t_{i+1}-t_i\le0.45/c^2$'' should be $\ge$.
\item Proof of Lemma~7.5 (p.~45): in the step for $B_{-(l+1)}$, ``Initially
  the count is at most $|B_{-(l+2)}|\le0.008|M|$'' should read
  $|B_{-(l+1)}|$.
\item Proof of Theorem~3.1 (p.~46): ``after Phase~7 and Phase~7'' should read
  ``after Phase~7 and Phase~8.''
\item Proof of Theorem~3.1 (p.~46): ``We finally justify that the expected
  stabilization time is $O(n\log n)$'' should read $O(\log n)$, as concluded
  at the end of the paragraph.
\item Very small populations: Phase~0 ends only through Clock--Clock
  interactions, and for $n=3$ (input $A,A,B$) it produces at most one Clock
  agent, so Phase~0 never ends. The sentence on p.~46, ``By Lemma~5.2, we
  must produce at least 2 Clock agents by the end of Phase~0,'' needs a lower
  bound on~$n$. This does not affect the asymptotic theorem.
\end{enumerate}
\end{erratum}

\section*{Formal verification}
\sloppy

The formalization is in the public repository
\url{https://github.com/zinan-huang/Ripple}, directory
\texttt{Ripple/PopulationProtocol/ExactMajority}. It follows the paper's
protocol and its proof phase by phase, with the corrections above and the
clock parameter $k=256$. The main theorem is
\texttt{TheoremThreeOne.exact\_majority\_theorem\_three\_one}. For every population
size $n$ above an explicit threshold it states, for one protocol: the number
of states is at most a constant times $\lceil\log_2 n\rceil+1$; the protocol
stably computes majority; and for every valid input of size $n$ the
stabilization time exceeds a constant times $n\log n$ interactions with
probability at most $110/n^2$, and its expectation is $O(n\log n)$
interactions, with explicit constants (not optimized). Correctness alone is
proved for all $n\ge11$. \texttt{\#print axioms} reports only
\texttt{propext}, \texttt{Classical.choice} and \texttt{Quot.sound}.

\end{document}
